%% =========================================================================
%% The Composite Floor: (C-infinity)
%% Promoted out of "Variants of the Conjecture" (2026-08-17 revision):
%% (C-infinity) is not a variant of MC, it is the necessary condition beneath
%% every route to it, and it carries the paper's longest unconditional arc.
%% =========================================================================
\section{The Composite Floor: $(\mathrm{C}\infty)$}
\label{sec:cinfty}
%% =========================================================================

\begin{roadmap}
Three independent lines of attack on Mullin's conjecture---the target
ladder, the smallness family, and the obstruction programme---terminate at
the same open statement:
\[
  (\mathrm{C}\infty):\qquad
  \Prod{2}(n)+1 \text{ is composite for infinitely many } n .
\]
This section is about that statement alone.  The arc:
\begin{enumerate}[nosep,leftmargin=*]
\item[\S\ref{sec:cinfty-beneath}] Everything rests on it, and the
  heuristics say it should be true with density one.
\item[\S\ref{sec:growth-floor}] What the smallness family actually needs is
  stronger: not that Euclid numbers are sometimes composite, but that they
  are sometimes \emph{smooth relative to their own size}.
\item[\S\ref{sec:defect-telescope}] Growth is governed by a single real
  number $C = \lim_N \log\Prod{2}(N)/2^N$, and $C = 0$ is \emph{equivalent}
  to $(\mathrm{C}\infty)$---so the failure branch is exactly the autonomous
  branch, and every obstruction in this paper bears on all of it.
\item[\S\ref{sec:sylvester-tower}] $(\mathrm{C}\infty)$ is Sylvester-tower
  primality: no orbit of $s \mapsto s^2-s+1$ through a Euclid number is
  all-prime.
\item[\S\ref{sec:backward-orbit}--\ref{sec:level-tower}] Restated in the
  paper's own language it is a \emph{hitting} statement for the residue
  walk, against a target of which the classical $\Phi_3$ death equation is
  only the first rung.  Rungs two and three are reachable by quadratic
  reciprocity; the ladder is finite.
\item[\S\ref{sec:arboreal}] The arboreal input---what Chebotarev would
  supply---is free.  The residual gap is disjoint from it.
\item[\S\ref{sec:ff-telescope}] Over $\FF_p[t]$, where every analytic input
  is a theorem, the whole telescope holds with the epsilons removed---and,
  uniformly in~$p$, nothing moves; for particular~$p$ it moves in both
  directions ($(\mathrm{C}\infty)$ false over $\FF_5[t]$, machine-checked;
  true over $\FF_2[t]$ and for $p \equiv 1 \pmod 3$).
\item[\S\ref{sec:cinfty-verdict}] The verdict: $\MC$ implies a
  no-infinite-prime-tower statement for a doubly exponential sequence,
  of the shape of the Fermat question.
\end{enumerate}
\end{roadmap}

%% =========================================================================
\subsection{Why Everything Rests on It}
\label{sec:cinfty-beneath}
%% =========================================================================

$(\mathrm{C}\infty)$ and its negation, perpetual primality, were defined in
\S\ref{sec:autonomous-branch}; the branch is not believed to hold, but no
current method excludes it.  Three separate routes to $\MC$ run into it.

\emph{The smallness family.}  Every statement of
\S\ref{sec:composite-floor}---$\mathrm{RD}$, $\mathrm{WM}$,
$\mathrm{MissingFinite}$, and $\MC$ itself---implies $(\mathrm{C}\infty)$,
and each of the three smallness statements fails outright on the branch.

\emph{The target ladder.}  Even (V)---the weakest orbit target of
\S\ref{sec:target-ladder}, asking only that every odd prime \emph{divide}
some Euclid number---forces $(\mathrm{C}\infty)$
(\lean{EM/Population/SylvesterTower.lean}{infinitelyManyComposite\_of\_everyPrimeDividesEuclid}).
On the perpetual branch, Bertrand supplies a prime strictly between
$\Prod{2}(T)+1$ and the next Euclid number: too large to divide the earlier
ones, too small to equal any later one, and those are all prime.

\emph{The obstruction programme.}  \S\ref{sec:anatomy-axis} shows
$(\mathrm{C}\infty)$ is what that programme leaves behind: it is not a
statement about the distribution of primes at all, but about the anatomy of
the numbers $\Prod{2}(n)+1$.

\smallskip\noindent
So the question is not whether $(\mathrm{C}\infty)$ is needed.  It is
whether it can be had.

\paragraph{Why a Euclid number should generally not be prime.}
The reason is not that primes are scarce among large integers.  It is
that this sequence \emph{penalises} primality, and at a doubly
exponential rate.  Compare the two branches of a single step.  If
$\Prod{2}(n)+1$ is composite, its least factor is a proper factor and the
accumulator grows by that factor.  If $\Prod{2}(n)+1$ is prime, the least
factor is the whole Euclid number, so
\[
  \Prod{2}(n+1) \;=\; \Prod{2}(n)\bigl(\Prod{2}(n)+1\bigr)
  \;>\; \Prod{2}(n)^2 ,
\]
and the accumulator \textbf{squares}
(\lean{EM/Population/CompositeFloor.lean}{sq\_lt\_prod\_succ\_of\_prime}).
Each prime Euclid number doubles $\log \Prod{2}(n)$, so the next
candidate---whose primality has probability of order $1/\log$---is twice
as unlikely.  Primality here does not merely occur rarely: it consumes
the very resource that makes it possible.

That is enough for an unconditional theorem.  Nothing need be assumed,
because the composite steps only help: they enlarge the accumulator too.

\begin{theorem}[Primality is self-limiting ---
\lean{EM/Population/CompositeFloor.lean}{two\_pow\_two\_pow\_primeEuclidCount\_le\_prod}]
\label{thm:self-limiting}
Let $m$ be the number of $n < N$ with $\Prod{2}(n)+1$ prime.  Then
\[
  2^{2^{m}} \;\leq\; \Prod{2}(N),
  \qquad\text{equivalently}\qquad
  m \;\leq\; \log_2 \log_2 \Prod{2}(N) .
\]
Prime Euclid numbers are at most \emph{doubly logarithmically} many in
the size of the accumulator.
\end{theorem}

\paragraph{The matching heuristic.}
It is worth seeing that this bound is essentially sharp, because the
naive sieve heuristic---done carefully---reproduces it.  The Euclid
number is coprime to every prime in the bag, so it is rough, and
roughness makes primality \emph{more} likely, not less: for a number of
size $X$ all of whose small prime factors have been excluded,
\[
  \mathbb{P}(\text{prime})
  \;\approx\;
  \frac{1}{\log X}\prod_{p \in S_n}\Bigl(1-\frac1p\Bigr)^{-1}.
\]
On the behaviour the conjecture predicts, the bag $S_n$ is essentially
the primes up to $p_n$, so by Mertens the correction is
$\asymp e^{\gamma}\log p_n$, while $\log \Prod{2}(n) = \theta(p_n) \asymp p_n$.
Hence
\[
  \mathbb{P}\bigl(\Prod{2}(n)+1 \text{ prime}\bigr)
  \;\approx\; \frac{e^{\gamma}\log p_n}{p_n}
  \;\asymp\; \frac{e^{\gamma}}{n},
\]
using $p_n \asymp n\log n$.  Two things follow, and they pull in
opposite directions.  The probability tends to~$0$, so the Euclid numbers
should be composite for \emph{almost every}~$n$---which is
$(\mathrm{C}\infty)$ with density one.  But $\sum 1/n$ diverges, so one
should expect \emph{infinitely many} prime Euclid numbers as well: the
sequence should enter the autonomous branch again and again, and leave it
every time.

(A different heuristic, recorded as Dead End~\#91, predicts only finitely
many prime Euclid numbers; it assumes doubly exponential growth
$\log \Prod{2}(n) \asymp 2^{n}$, i.e.\ the perpetual-primality branch,
under which $\sum 1/\log\Prod{2}(n)$ converges.  The two are not in
conflict: they condition on the two branches of the dichotomy of
\S\ref{sec:sylvester-tower}, and on the branch MC predicts the count
diverges.)  The two estimates agree.  The heuristic predicts $\asymp \log N$ primes
among the first $N$ Euclid numbers; Theorem~\ref{thm:self-limiting} caps
that count by $\log_2\log_2 \Prod{2}(N) \asymp \log_2(n \log n) \asymp
\log N$ on the same behaviour.  So the unconditional bound is saturated
by the expected truth, and cannot be improved without genuinely new
input---which is the same conclusion, reached from the opposite side, as
\S\ref{sec:sylvester-tower}.

\paragraph{Can either half be proved?}
Neither.  But they are not symmetric, and it is worth being precise about
how.  The prime and composite stages partition~$\NN$, so
\emph{at least one} of ``infinitely many prime Euclid numbers'' and
$(\mathrm{C}\infty)$ is true---that much is free.  Only one of the two
disjuncts carries content: if $(\mathrm{C}\infty)$ fails, MC fails and
every smallness statement of \S\ref{sec:composite-floor} fails with it; if instead only
finitely many Euclid numbers are prime, nothing at all follows.  The
free disjunction therefore says: \emph{either the interesting statement
holds, or the harmless one does}, which is no help.

For the prime half there is no technique to try.  Every method that
produces infinitely many primes in a sparse sequence needs one of two
things: a positive-density family to sieve, or an algebraic structure
supporting a Chebotarev-type equidistribution.  The Euclid numbers offer
neither---there is one integer per index, and the recursion defining them
invokes factorisation itself.  Fermat, Mersenne and primorial numbers are
far more structured, and the corresponding questions are open for all
three.

\smallskip\noindent
So $(\mathrm{C}\infty)$ is required, expected to hold with density one, and
carries the whole content of the disjunction.  The rest of this section
takes it apart.  The first surprise is that the smallness family does not
need it: it needs something stronger.

%% =========================================================================
\subsection{The Floor Is a Growth Statement}
\label{sec:growth-floor}
%% =========================================================================

The floor of \S\ref{sec:composite-floor} was proved by showing that on the
perpetual-primality branch $\sum_k 1/\seq{2}(k)$ converges.  It is worth
asking what that argument actually consumed, because it is not primality.
Perpetual primality enters only through the inequality
$\seq{2}(n+1) \geq 2^{\,n}$.  Any lower bound on the selected primes
would do, and ``large'' may be measured against \emph{any} benchmark
whose reciprocals converge.  Stripping the hypothesis to what the proof
uses gives the following, whose proof is a comparison of series with no
arithmetic in it at all.

\begin{keyresult}
\textbf{The growth floor} $(\mathrm{S})$
(\lean{EM/Population/CompositeFloor.lean}{growth\_floor\_landscape}).
Let $f$ be positive with $\sum_n 1/f(n) < \infty$.  If
$f(n) \leq \seq{2}(n)$ for all large $n$, then
$\sum_k 1/\seq{2}(k) < \infty$
(\lean{EM/Population/CompositeFloor.lean}{summable\_one\_div\_seq\_of\_lower\_bound}).
Contrapositively, $\mathrm{RD}$ forces the selected primes below every
such benchmark infinitely often
(\lean{EM/Population/CompositeFloor.lean}{exists\_lt\_of\_reciprocalDivergence}); in
arithmetic terms, for every fixed $c$,
\[
  \minFac\bigl(\Prod{2}(n)+1\bigr) \;<\; 2^{\,n-c}
  \qquad\text{for infinitely many } n
\]
(\lean{EM/Population/CompositeFloor.lean}{exists\_small\_minFac\_of\_reciprocalDivergence}).
\end{keyresult}

\noindent
This is strictly stronger than $(\mathrm{C}\infty)$, which asks only that
the least factor be \emph{proper}.  Indeed $(\mathrm{S})$ with $c = 0$
returns $(\mathrm{C}\infty)$ at once, since $2^{\,n} < \Prod{2}(n)+1$
(\lean{EM/Population/CompositeFloor.lean}{infinitelyManyComposite\_of\_reciprocalDivergence\_via\_growth}),
and perpetual primality is recovered as the extreme case
$f(n) = 2^{\,n}$.  So the true content of the floor is not that the
Euclid numbers are sometimes composite, but that they are sometimes
\emph{smooth relative to their own size}.  That is a strictly harder
demand, and it is what the whole smallness family rests on.

%% =========================================================================
\subsection{The Defect Telescope}
\label{sec:defect-telescope}
%% =========================================================================

\paragraph{The criterion, and why it looked lossy.}
The self-limiting bound of \S\ref{sec:cinfty-beneath} yields a genuine
reformulation of $(\mathrm{C}\infty)$.  Splitting the first $N$ stages into
prime and composite ones and applying Theorem~\ref{thm:self-limiting},
\[
  N \;\leq\; \#\{n < N : \Prod{2}(n)+1 \text{ composite}\}
      \;+\; \log_2\log_2 \Prod{2}(N)
\]
(\lean{EM/Population/CompositeFloor.lean}{le\_compositeEuclidCount\_add\_log\_log}),
so composite Euclid numbers occur past every stage as soon as $N$ outruns
$\log_2\log_2 \Prod{2}(N)$ by an unbounded margin
(\lean{EM/Population/CompositeFloor.lean}{infinitelyManyComposite\_of\_subtower\_growth}).
\emph{$(\mathrm{C}\infty)$ follows from any bound saying the accumulator
grows strictly slower than the maximal tower rate.}

This reformulation is honest about its own difficulty.  The selected
factor never exceeds the Euclid number, so the accumulator obeys the
unconditional ceiling
\[
  \Prod{2}(N) + 1 \;\leq\; 3^{\,2^{N}}
\]
(\lean{EM/Population/CompositeFloor.lean}{prod\_add\_one\_le\_three\_pow}),
whence $\log_2\log_2 \Prod{2}(N) \leq N + O(1)$ and the displayed
inequality is vacuous unless the ceiling can be beaten by an unbounded
margin.  Beating it is precisely what perpetual primality forbids, so the
reformulation ought to be equivalent in strength to $(\mathrm{C}\infty)$
rather than weaker.  The rest of this subsection proves that it is,
exactly, and in the process shows that the change of subject---from the
primality of individual terms to the growth rate of one explicit
recursion---is information-preserving rather than lossy.

\smallskip\noindent
What follows makes the growth side into an object.  It attaches to the
orbit a single real number, shows that the criterion above is
\emph{exactly} the vanishing of that number, and---this is the point---%
shows that the number is forced to vanish by hypotheses far weaker than
$\MC$.  Everything here is elementary; none of it is about the primality of
any individual term.

\paragraph{The identity.}
Write $L(n) = \log \Prod{2}(n)$.  Since
$\Prod{2}(n+1) = \Prod{2}(n)\,\seq{2}(n+1)$,
\[
  L(n+1) \;=\; L(n) + \log \seq{2}(n+1).
\]
The maximal step is the perpetual-primality step
$\seq{2}(n+1) = \Prod{2}(n)+1$, where the second summand is $L(n)$ again
and $L$ doubles.  Measuring every step against that maximum defines the
\textbf{defect}
\[
  \delta(n) \;:=\; L(n) - \log \seq{2}(n+1),
\]
and the recursion becomes exactly
\[
  L(n+1) \;=\; 2L(n) - \delta(n)
\]
(\lean{EM/Population/DefectTelescope.lean}{logProd\_succ\_eq\_two\_mul\_sub\_defect}).
Dividing by $2^{\,n+1}$ telescopes it:
\begin{equation}
\label{eq:defect-telescope}
  \frac{L(N)}{2^{N}}
  \;=\; L(0) \;-\; \sum_{n<N} \frac{\delta(n)}{2^{\,n+1}}
\end{equation}
(\lean{EM/Population/DefectTelescope.lean}{normLog\_eq\_sub\_sum}).
The defect is not quite nonnegative---a prime Euclid number overshoots,
because $\seq{2}(n+1) = \Prod{2}(n)+1 > \Prod{2}(n)$---but it overshoots
by at most $\log(1 + 1/\Prod{2}(n)) \leq 2^{-(n+1)}$, using the
unconditional bound $2^{\,n+1} \leq \Prod{2}(n)$
(\lean{EM/Population/DefectTelescope.lean}{neg\_two\_pow\_le\_defect}).  So
$L(N)/2^{N}$ is decreasing up to a summable error, is positive, and
therefore \textbf{converges}.

\begin{definition}[The growth constant]
\label{def:growth-constant}
$C \;:=\; \displaystyle\lim_{N\to\infty} \frac{\log \Prod{2}(N)}{2^{N}}
 \;=\; L(0) - \sum_{n\geq 0} 2^{-(n+1)}\delta(n) \;\geq\; 0$
(\lean{EM/Population/DefectTelescope.lean}{tendsto\_normLog},
 \lean{EM/Population/DefectTelescope.lean}{growthConstant\_nonneg}).
\end{definition}
\noindent The existence of such a limit for a doubly exponential
recurrence $x_{n+1} = x_n^2(1+o(1))$ is the standard Aho--Sloane
argument~\cite{AhoSloane1973}; what is specific here is the exact
defect telescope and the identification of $C = 0$ below.

\noindent
Note the weighting.  A defect incurred at stage $n$ is discounted by
$2^{-n}$, so \emph{no finite computation can contribute to $C$}: the
first fifty stages move it by an amount that the tail can undo.  This is
not a defect of the definition but a correct reflection of the problem,
and it is the reason the numerical evidence of \S\ref{sec:hard} is
silent here.

\paragraph{What the constant decides.}
Three statements, all proved, and all about $C$ alone.

\begin{keyresult}
\textbf{The defect landscape}
(\lean{EM/Population/DefectTelescope.lean}{defect\_landscape}).
\begin{enumerate}[nosep,leftmargin=*]
\item The sub-tower growth criterion of \S\ref{sec:composite-floor}---%
that $N$ outruns $\log_2\log_2\Prod{2}(N)$ by an unbounded margin---holds
\emph{if and only if} $C = 0$
(\lean{EM/Population/DefectTelescope.lean}{subtower\_growth\_iff\_growthConstant\_eq\_zero}).
\item $C = 0 \Rightarrow (\mathrm{C}\infty)$
(\lean{EM/Population/DefectTelescope.lean}{infinitelyManyComposite\_of\_growthConstant\_eq\_zero}).
\item $\mathrm{RD} \Rightarrow C = 0$
(\lean{EM/Population/DefectTelescope.lean}{growthConstant\_eq\_zero\_of\_reciprocalDivergence}),
hence so do $\mathrm{WM}$, $\mathrm{MissingFinite}$ and $\MC$.
\end{enumerate}
\end{keyresult}

\paragraph{The adelic reading.}
The greedy step multiplies by the principal idele $p_{n+1}$: it dilates
the archimedean coordinate by $p_{n+1}$ and contracts the $p_{n+1}$-adic
one by the same factor, the product formula balancing the two.  In that
language the defect is, up to $\log(1 + 1/\Prod{2}(n))$, the logarithm of
the \emph{discarded cofactor} $(\Prod{2}(n)+1)/p_{n+1}$
(\lean{EM/Population/AdelicShadow.lean}{defect\_eq\_log\_cofactor}): the
archimedean size of the part of the divisor of the Euclid number that the
rule throws away.  The telescope is the product formula summed along the
orbit with the chosen place removed at each step, and $C$ measures
discard, not location --- which is why the growth axis is a complete
invariant of perpetual primality and of nothing finer.  On the local side
the same bookkeeping gives, for every seed $m \ge 1$ and prime~$q$, that
the Euclid numbers are divisible by~$q$ at only finitely many steps
(\lean{EM/Population/AdelicShadow.lean}{hits\_finite}; the multipliers at
the hits are distinct primes $\le q$), that more than $\pi(q-1)$ hits force
the selection of~$q$
(\lean{EM/Population/AdelicShadow.lean}{captured\_of\_many\_hits}), and
hence that the failure of $\MC$ at~$q$ is exactly ``$\Prod{}(n)$ and,
eventually, $\Prod{}(n)+1$ are $q$-adic units''
(\lean{EM/Population/AdelicShadow.lean}{misses\_iff\_eventually\_unit\_both}).
This is the cleanest orbit-level form of the failure of $\MC$ we have; it
is also symmetric between capture and failure, which is the precise sense
in which the product formula localises nothing.

\noindent
Item (1) is what the previous subsection could only gesture at.  The
criterion that looked merely sufficient for $(\mathrm{C}\infty)$ is
\emph{identical} to $C = 0$; and if $C > 0$ then
$\log_2\log_2 \Prod{2}(N) = N + \log_2 C + \log_2\log_2 e + o(1)$ (natural
$\log$ in $C$), which is exactly why the
inequality $N \leq \#\{\text{composites}\} + \log_2\log_2\Prod{2}(N)$ is
vacuous on that branch.

Item (3) is the new floor, and it is strictly below $(\mathrm{S})$.  Where
$(\mathrm{S})$ says the selected primes are sometimes small,
$\mathrm{RD} \Rightarrow C = 0$ says $\log \Prod{2}(N) = o(2^{N})$: a
statement about the orbit as a whole, with no reference to any individual
term.  The chain
\[
  \MC \;\Rightarrow\; \mathrm{RD} \;\Rightarrow\; C = 0
  \;\Rightarrow\; (\mathrm{C}\infty)
\]
replaces an anatomy question by a growth question at both ends.

\paragraph{How wide is the failure branch?  Not at all.}
The natural next question is what the surviving alternative looks like.
If $C > 0$ then the terms $2^{-(n+1)}\delta(n)$ of
\eqref{eq:defect-telescope} are summable, so
$\delta(n) = o(2^{n}) = o(L(n))$, that is
\begin{equation}
\label{eq:near-perpetual}
  \frac{\log \seq{2}(n+1)}{\log \Prod{2}(n)} \;\longrightarrow\; 1
\end{equation}
(\lean{EM/Population/DefectTelescope.lean}{tendsto\_log\_seq\_div\_logProd\_of\_pos}).
Read naively this looks like a branch far wider than perpetual primality:
the least prime factor need only be $\Prod{2}(n)^{1-o(1)}$, not the Euclid
number itself, and on such a branch no Euclid number need be prime.

That reading is wrong, and the fact that kills it is the oldest one in the
subject.  \emph{Trial division}: if $\minFac(X) > \sqrt X$ then $X$ is
prime, since otherwise $X = pm$ with $m \geq p > \sqrt X$ forces $X > X$
(\lean{EM/Core/Euclid.lean}{minFac\_sq\_le}).  A ratio in
\eqref{eq:near-perpetual} eventually above $1/2$ therefore \emph{is}
primality.  Quantitatively, the defect has a gap.

\begin{lemma}[The defect gap ---
\lean{EM/Population/DefectTelescope.lean}{defect\_gap}]
\label{lem:defect-gap}
For every $n$, either $\delta(n) \leq 0$---which happens exactly when
$\Prod{2}(n)+1$ is prime---or
\[
  \delta(n) \;\geq\; \tfrac12 L(n) - 2^{-(n+2)} ,
\]
which happens exactly when it is composite.  Nothing lies in between: the
ratio \eqref{eq:near-perpetual} never enters $(\tfrac12 + \varepsilon,\, 1)$.
\end{lemma}

\noindent
Feeding the composite case back into the telescope, one composite stage
multiplies the (corrected) normalised accumulator by $3/4$.  Iterating
gives an unconditional bound in which composite stages \emph{provably damp
the tower}:
\begin{equation}
\label{eq:damping}
  \log \Prod{2}(N)
  \;\leq\;
  \Bigl(\log 2 + \tfrac13\Bigr)\,
  \Bigl(\tfrac34\Bigr)^{\#\{n<N\,:\,\Prod{2}(n)+1 \text{ composite}\}}
  \, 2^{N}
\end{equation}
(\lean{EM/Population/DefectTelescope.lean}{logProd\_le\_pow\_mul}).  This
refines the tower ceiling $\Prod{2}(N)+1 \leq 3^{2^N}$ of
\S\ref{sec:composite-floor}, which is the case of no composite stages at
all.  And it closes the circle: infinitely many composite stages drive the
right-hand side of \eqref{eq:damping} to $o(2^N)$, so $C = 0$.

\begin{keyresult}
\textbf{The growth constant is a complete invariant}
(\lean{EM/Population/DefectTelescope.lean}{infinitelyManyComposite\_iff\_growthConstant\_eq\_zero},
 \lean{EM/Population/DefectTelescope.lean}{growthConstant\_pos\_iff\_perpetualPrimality}).
\[
  C = 0 \iff (\mathrm{C}\infty),
  \qquad
  C > 0 \iff \Prod{2}(n)+1 \text{ is prime for all large } n .
\]
\end{keyresult}

\begin{theorem}[The dichotomy, sharp ---
\lean{EM/Population/DefectTelescope.lean}{defect\_dichotomy}]
\label{thm:defect-dichotomy}
Either $(\mathrm{C}\infty)$ holds and $C = 0$; or the sequence is
eventually on the perpetual-primality branch, $C > 0$, and
$\minFac(\Prod{2}(n)+1) = \Prod{2}(n)^{1-o(1)}$.  Each alternative is
equivalent to its growth-side description.
\end{theorem}

\paragraph{What this settles, and what it costs.}
It costs the hope that the growth reformulation was weaker than
$(\mathrm{C}\infty)$.  It is not weaker by any margin: $C$ is a complete
invariant, so \S\ref{sec:composite-floor}'s remark that the reformulation
is ``equivalent in strength'' is now literally a theorem.  Together with
\S\ref{sec:target-ladder} and \S\ref{sec:composite-floor} that is three
attempted weakenings, all of which came back exactly equivalent, for the
same structural reason: the problem has very few genuine degrees of
freedom.

What it settles is more useful.  The failure branch is \emph{exactly} the
autonomous branch.  Every obstruction in this paper---the autonomous-map
argument of \S\ref{sec:hard}, the density-$1/2$ $\Phi_3$ exclusion, the
$\Phi_6$ orbit counting of \S\ref{sec:sylvester-tower}, the refutation
below stage~$8$---applies only where the recursion degenerates into an
autonomous map $w \mapsto w^2+w$, and one might have feared that this was
a measure-zero corner of the true failure set.  It is the whole of it.
Nothing escapes, and the machinery is aimed correctly.

Finally, the route to item (3) is worth isolating, because it converts a
growth statement back into an arithmetic one without using trial division
at all.  On the branch \eqref{eq:near-perpetual} the selected primes
eventually exceed $(\sqrt 2)^{\,n}$---geometric growth, with room to
spare---and $(\mathrm{S})$ then makes $\sum_k 1/\seq{2}(k)$ converge
(\lean{EM/Population/DefectTelescope.lean}{not\_reciprocalDivergence\_of\_growthConstant\_pos}).
Given the equivalence this is a second proof of $\mathrm{RD} \Rightarrow
(\mathrm{C}\infty)$, and the only one that is insensitive to how large the
least prime factor is allowed to be.

\paragraph{The seeded growth constant: $C$ as a factor map.}
\label{sec:seeded-growth}
Nothing in the telescope uses the seed.  For every $m \ge 2$ the same
argument gives $C(m) = \lim_n \log(T^n m)/2^n$ with $T(n) = n \cdot
\minFac(n{+}1)$ (\code{EM/Population/SeededGrowth.lean}, \code{sgrowth}),
and the tail identity $T^{n+1} m = T^n(Tm)$ turns the limit into a
\emph{semiconjugacy}: $C(Tm) = 2\,C(m)$ and $C(T^k m) = 2^k C(m)$
(\lean{EM/Population/SeededGrowth.lean}{sgrowth\_T},
 \lean{EM/Population/SeededGrowth.lean}{sgrowth\_iterate}).  The complete
invariance holds seed by seed---$C(m) = 0$ iff the orbit of $m$ has
infinitely many composite Euclid numbers, $C(m) > 0$ iff it is eventually
on the perpetual-primality branch
(\lean{EM/Population/SeededGrowth.lean}{seedInfinitelyManyComposite\_iff\_sgrowth\_eq\_zero},
 \lean{EM/Population/SeededGrowth.lean}{sgrowth\_pos\_iff\_eventually\_prime})---and
$C(2)$ is the constant above
(\lean{EM/Population/SeededGrowth.lean}{sgrowth\_two}).  So $C$ maps
$(X,T)$ onto doubling on $[0,\infty)$; the set $\{C > 0\}$ is
$T$-invariant, and \defname{MixedDiversity} (\S\ref{sec:variants}), the
statement that $(\mathrm{C}\infty)$ holds from \emph{every} seed, is
exactly the statement that this invariant set is empty
(\lean{EM/Population/SeededGrowth.lean}{mixedDiversity\_iff\_sgrowth\_zero}).
One consequence, and one non-consequence.  A seed with $C(m) > 0$ would be a doubly exponential
Sylvester-type prime tower; its non-existence for every seed is a
universal Fermat-shaped statement, strictly stronger than the floor
$(\mathrm{C}\infty)$ for the seed~$2$ and independent of MC.  (We do
\emph{not} claim an ergodic content: $(X,T)$ has no invariant probability
measure at all, trivially, because $T(m) > m$ on a countable set; the
semiconjugacy to doubling adds nothing to that.)

\paragraph{How thin is $\{C>0\}$?}
\label{sec:growth-density}
Threshold by threshold, null.  Writing $E_N = \{m : T^n m + 1 \text{
prime for all } n \ge N\}$, one has $\{C > 0\} = \bigcup_N E_N$
(\lean{EM/Population/GrowthDensity.lean}{sgrowth\_pos\_subset\_iUnion}),
and every $E_N$ has natural density zero
(\lean{EM/Population/GrowthDensity.lean}{hasDensityZero\_perpetual}).
The proof uses no prime number theorem: primes have density zero by the
elementary sieve bound with the primorial modulus and $\prod_{r<z}(1-1/r)
\to 0$ (\lean{EM/Population/GrowthDensity.lean}{hasDensityZero\_prime}),
and density zero pulls back under $T$
(\lean{EM/Population/GrowthDensity.lean}{hasDensityZero\_comp\_T}: split by
$p = \minFac(m+1)$; each fixed $p \le z$ injects $m \mapsto mp$ into a thin
set, and the tail $\minFac(m+1) > z$ is the rough count of density
$c(z{+}1) \to 0$), so at every fixed stage almost no seed has a prime
Euclid number
(\lean{EM/Population/GrowthDensity.lean}{hasDensityZero\_genProd\_prime}).
What does \emph{not} follow, and is not claimed: that $\{C>0\}$ itself is
null.  Natural density is not countably subadditive, and the pull-back
costs a factor depending on the head cut $z$ that iterating $N$ times does
not control uniformly.  ``MixedDiversity holds for almost every seed'' is
therefore a threshold-by-threshold statement here, and the uniform version
is a genuine open question, though a mild one next to MC.

\paragraph{What the growth coordinate says about the walk.}
\label{sec:size-residue}
The two projections are coupled through exactly one bit.  At a
\emph{prime} stage the multiplier is $P+1$ itself, so its residue is forced
by the walk position, $\minFac(P+1) \equiv w+1$
(\lean{EM/Population/SizeResidueDecoupling.lean}{multiplier\_residue\_of\_prime\_stage}):
this is the autonomous branch $w \mapsto w(w+1)$ of
\S\ref{sec:autonomous-branch}, the branch on which $C>0$ lives.  At a
\emph{composite} stage nothing is forced.  In the ensemble, for every odd
prime~$q$, every walk position $w \ne -1$, every unit residue~$a$ and every
bound~$K$ there is a squarefree seed $m \equiv w \pmod q$ whose Euclid
number is composite with $\minFac(m+1) \equiv a \pmod q$ and
$\minFac(m+1) \ge K$
(\lean{EM/Population/SizeResidueDecoupling.lean}{exists\_seed\_composite\_residue\_size};
take $m = 2m'$ with $m'$ a Dirichlet prime in a CRT class that makes
$2m'+1 \equiv 0$ modulo a Dirichlet prime $p \equiv a$, $p > K$, and
$\equiv 2$ modulo every other odd prime below~$p$).  So on $\{C = 0\}$ the
growth projection carries no residue information beyond ``the walk is not on
its autonomous branch at these stages''---which is why the two projections
are not on a par, and why nothing above the floor came from the growth
side.  This is a statement about seeds; along the orbit of~$2$ the pairs
(position, size) are what they are (Dead Ends~\#90/\#117).

\paragraph{A second invariant: relative size.}
\label{sec:relative-size}
$C$ is scaled by~$T$.  An invariant that is \emph{not} scaled is the
asymptotic relative size of the selected multiplier,
\[
  \rho(m) \;=\; \liminf_{n}\;
  \frac{\log \minFac(T^n m + 1)}{\log T^n m},
\]
for which $\rho(Tm) = \rho(m)$ exactly, since $T$ acts on the orbit as the
shift (\lean{EM/Population/RelativeSize.lean}{rho\_T}).  It refines the
dichotomy: $C(m) > 0 \iff \rho(m) = 1$ (on the perpetual-primality branch
the ratio tends to~$1$), $C(m) = 0 \iff \rho(m) \le \tfrac12$ (at a
composite stage $\minFac(P+1)^2 \le P+1$), so $\rho \in [0,\tfrac12] \cup
\{1\}$
(\lean{EM/Population/RelativeSize.lean}{sgrowth\_pos\_iff\_rho\_eq\_one},
 \lean{EM/Population/RelativeSize.lean}{sgrowth\_eq\_zero\_iff\_rho\_le\_half},
 \lean{EM/Population/RelativeSize.lean}{rho\_dichotomy}); and reciprocal
divergence forces $\rho = 0$, because $\rho > \varepsilon$ makes
$\minFac(T^n m+1) \ge (T^n m)^{\varepsilon} \ge 2^{\varepsilon n}$
eventually and the reciprocals summable
(\lean{EM/Population/RelativeSize.lean}{rho\_eq\_zero\_of\_seedRD}).  For
the orbit of~$2$ this inserts a rung in the floor ladder
(\lean{EM/Population/RelativeSize.lean}{relative\_size\_landscape}):
\[
  \MC \;\Rightarrow\; \mathrm{RD} \;\Rightarrow\; \rho(2) = 0
  \;\Rightarrow\; \rho(2) \le \tfrac12 \;\iff\; (\mathrm{C}\infty)
  \;\iff\; C = 0 .
\]
Whether $\rho(2) = 0$ is strictly weaker than RD, or strictly stronger
than $(\mathrm{C}\infty)$, is open; (S) sits between RD and $\rho(2)=0$.
Nothing here points upward: $\rho$ is a growth-side quantity and inherits
the decoupling of the previous paragraph.

%% =========================================================================
\subsection{$(\mathrm{C}\infty)$ Identified: the Sylvester Tower}
\label{sec:sylvester-tower}
%% =========================================================================

The floor is worth naming, because it is not new.  On the
perpetual-primality branch the Euclid numbers satisfy
\[
  \bigl(\Prod{2}(n+1)+1\bigr)
  = \bigl(\Prod{2}(n)+1\bigr)^{2} - \bigl(\Prod{2}(n)+1\bigr) + 1,
\]
which is \emph{Sylvester's recursion} $s \mapsto s^2 - s + 1$.  And
$\Prod{2}(n)+1$ for $n = 0,1,2,3$ is
\[
  3,\quad 7,\quad 43,\quad 1807,
\]
which is Sylvester's sequence $2, 3, 7, 43, 1807, 3263443, \dots$ from
its second term.  The Euclid--Mullin sequence \emph{is} Sylvester's
sequence for exactly as long as the Euclid numbers stay prime, and it
broke away at $1807 = 13 \cdot 139$---which is why $\seq{2}(4) = 13$.

\begin{theorem}[(C$\infty$) is Sylvester-tower primality ---
\lean{EM/Population/SylvesterTower.lean}{infinitelyManyComposite\_iff\_tower\_composite}]
\label{thm:sylvester-tower}
Perpetual primality from~$N$ holds if and only if every term of the
Sylvester tower seeded at $\Prod{2}(N)+1$ is prime.  Hence
(C$\infty$) holds if and only if, for every~$N$, the tower seeded at
$\Prod{2}(N)+1$ contains a composite term.
\end{theorem}

\noindent
This closes a loop with the take-all rule of the previous subsection,
which \emph{is} Sylvester's sequence: the $\minFac$ rule degenerates to
the take-all rule precisely on the perpetual-primality branch, which is
the one branch where the min rule's growth advantage evaporates.
Whether Sylvester's own sequence contains infinitely many primes (or
only finitely many, or is even squarefree) is open~\cite{Vardi1991};
what is known is that the primes dividing some Sylvester number have
density zero~\cite{Odoni1985}.

\paragraph{Why the two elementary attacks are unavailable.}
Both standard routes to a compositeness theorem are provably blocked
here, and by obstructions this paper has already established.

\emph{Congruence.}  Forcing a small prime $\ell$ into $\Prod{2}(n)+1$
means steering $\Prod{2}(n) \bmod \ell$ into a prescribed set of residues.
Taken naively that set is the root set of $\Phi_6$, non-empty only for
$\ell \equiv 1 \pmod 6$---and \S\ref{sec:backward-orbit} shows the real
target is considerably larger than that.  But enlarging it does not help,
because the obstruction is not the size of the target: it is that the
requirement is a statement about the position of \emph{one specific orbit}
modulo a prime outside the bag, and
\lean{EM/Obstruction/NoInvariant.lean}{free\_transition} says congruence
data does not constrain it.  \S\ref{sec:backward-orbit}--\ref{sec:arboreal}
make that verdict progressively sharper; none of them overturns it.

\emph{Reciprocity.}  The available symbol data is automatically
consistent.  Since $\Prod{2}(n) \equiv 2 \pmod 4$, the candidate
$p = \Prod{2}(n)+1$ satisfies $p \equiv 3 \pmod 4$ and $(p-1)/2$ is odd;
since $p \equiv 1 \pmod{\seq{2}(k)}$ for every $k \leq n$, quadratic
reciprocity gives $\bigl(\tfrac{\seq{2}(k)}{p}\bigr) = (-1)^{(\seq{2}(k)-1)/2}$
and hence
\[
  \Bigl(\tfrac{\Prod{2}(n)}{p}\Bigr) = \Bigl(\tfrac{2}{p}\Bigr) \cdot (-1)^{t},
  \qquad t = \#\{k \leq n : \seq{2}(k) \equiv 3 \bmod 4\}.
\]
But $\bigl(\tfrac{\Prod{2}(n)}{p}\bigr) = \bigl(\tfrac{-1}{p}\bigr) = -1$,
and $\bigl(\tfrac{2}{p}\bigr)$ is fixed by $\Prod{2}(n) \bmod 8$, which is
fixed by the \emph{same} parity~$t$.  The two sides agree identically.
This is the concrete instance of
\lean{EM/Reciprocity/NoReciprocityInvariant.lean}{no\_reciprocity\_induction\_proof}.

%% =========================================================================
\subsection{$(\mathrm{C}\infty)$ as a Hitting Statement}
\label{sec:backward-orbit}
%% =========================================================================

The congruence attack is blocked, but the target it aims at is far smaller
than the one actually available, and it is worth saying by how much.  A
prime $\ell$ divides the $k$-th term of the tower seeded
at $s$ exactly when $s \bmod \ell$ reaches $0$ after $k$ steps of
$\Phi_6$, because reduction mod $\ell$ commutes with the recursion
(\lean{EM/Population/BackwardOrbit.lean}{cast\_tower},
 \lean{EM/Population/BackwardOrbit.lean}{dvd\_tower\_iff}).  Write
\[
  \mathrm{PreZero}(\ell)
  \;=\;
  \{x \in \ZZ/\ell\ZZ \;:\; \Phi_6^{\,k}(x) = 0 \text{ for some } k \geq 0\}
\]
for the \emph{backward orbit of~$0$}.  Then (C$\infty$) is exactly a
hitting statement for the residue walk.

\begin{theorem}[(C$\infty$) is walk-hitting ---
\lean{EM/Population/BackwardOrbit.lean}{backwardOrbitHitting\_iff\_infinitelyManyComposite}]
\label{thm:backward-orbit}
$(\mathrm{C}\infty)$ holds if and only if for every~$N$ there is a
prime~$\ell$ and a level~$k$ with
$\Phi_6^{\,k}\bigl(\walkZ{2}{\ell}{N}+1\bigr) = 0$, witnessed below the tower
term it kills.  Contrapositively, on the perpetual-primality branch the
walk value at the branch point satisfies
\[
  \walkZ{2}{\ell}{N} + 1 \;\notin\; \mathrm{PreZero}(\ell)
  \qquad\text{for \emph{every} prime } \ell \leq \Prod{2}(N)
\]
(\lean{EM/Population/BackwardOrbit.lean}{walkZ\_notMem\_preZero\_of\_perpetual}).
\end{theorem}

\noindent
The equivalence is exact, so nothing is lost; what changes is the shape
of what must be proved, and the comparison is stark.

\smallskip
\begin{center}
\begin{tabular}{@{}l@{\quad}l@{\quad}l@{}}
\toprule
\textbf{Statement} & \textbf{Target in $\ZZ/q\ZZ$} & \textbf{Moduli needed} \\
\midrule
$\MC$ & the single residue $-1$ & \emph{every} prime $q$ \\
$(\mathrm{C}\infty)$ & $\mathrm{PreZero}(\ell)$ & \emph{some} prime $\ell$, per stage \\
\bottomrule
\end{tabular}
\end{center}

\smallskip\noindent
Two features make the second target the larger one.  First,
$\mathrm{PreZero}$ is backward-closed under $\Phi_6$
(\lean{EM/Population/BackwardOrbit.lean}{mem\_preZero\_of\_phi6\_mem}): it
is the union of all iterated preimages of $0$, not a single residue.
Second---and this locates the classical obstruction precisely---%
\[
  \Phi_6(w+1) \;=\; w^2 + w + 1 \;=\; \Phi_3(w)
\]
(\lean{EM/Population/BackwardOrbit.lean}{phi6\_add\_one}), so \emph{level
one} of $\mathrm{PreZero}$ is exactly the death equation of the take-all
rule, the one whose insolubility modulo $q \equiv 2 \pmod 3$ drives the
density-$1/2$ failure of \S\ref{sec:selection-landscape}.  Every argument
the project has aimed at $\Phi_3$ uses one level; $(\mathrm{C}\infty)$ may
use any level.

Against that, the supply of usable moduli is halved: a nonzero element of
$\mathrm{PreZero}(\ell)$ forces a root of $\Phi_6$, hence a primitive
sixth root of unity, hence $6 \mid \ell-1$ (for primes $\ell \ne 2, 3$;
\lean{EM/Population/BackwardOrbit.lean}{six\_dvd\_sub\_one\_of\_phi6\_root}).
For $\ell \not\equiv 1 \pmod 6$ the target is $\{0\}$ and useless.  Even
so, heuristically the in-tree of a node in the functional graph of a map
on $\ZZ/\ell\ZZ$ has size of order $\sqrt{\ell}$, so the target density is
about $\ell^{-1/2}$ rather than $\ell^{-1}$, and
$\sum_{\ell \equiv 1 (6)} \ell^{-1/2}$ diverges---at rate
$\sqrt L/\log L$ rather than $\log\log L$.  (That count is a heuristic and
is not formalised; the criterion, the backward-closure and the $\Phi_3$
identification are.)

The conclusion is a calibration, not a proof.  $(\mathrm{C}\infty)$ needs
an \emph{existential} hit against a large target; $\MC$ needs a
\emph{universal} hit against a single residue.  Every hitting hypothesis
in this paper is therefore far more than $(\mathrm{C}\infty)$ requires---%
and none of them is available, because the obstruction is not the size of
the target but the fact that one specific orbit is being asked about.

%% =========================================================================
\subsection{The Level Tower: Past $\Phi_3$}
\label{sec:level-tower}
%% =========================================================================

\paragraph{One map, arrived at three times.}
Before climbing, it is worth pausing on a coincidence that this paper has
so far kept hidden by naming the same object differently in three places.
The map
\[
  q(w) \;=\; w^2 + w
\]
has appeared as: the \emph{residue walk on the perpetual-primality branch},
where the recursion degenerates and $\walkZ{2}{q}{n+1} = f(\walkZ{2}{q}{n})$
with $f(w) = w(w+1)$ (\S\ref{sec:autonomous-branch}); the walk of the
\emph{take-all selection rule}, which banks every prime factor and closes
the recursion into the same autonomous map
(\S\ref{sec:selection-landscape}); and now, conjugated by $w \mapsto w+1$,
the map whose backward orbit of $-1$ is the $(\mathrm{C}\infty)$ target.
They are the same map, and its death condition $q(w) = -1$, i.e.\
$\Phi_3(w) = 0$, is the same equation in all three readings.  From here on
we write $q$ for it throughout.

\paragraph{Raising the level past $\Phi_3$.}
The comparison above is only useful if the higher levels are genuinely
occupied, and they can be located exactly.  Translation by~$1$ conjugates
$q$ to $\Phi_6$, since $\Phi_6(w+1) = (w+1)w + 1 = q(w)+1$; iterating,
$\Phi_6^{\,k}(w+1) = q^{\,k}(w) + 1$
(\lean{EM/Population/BackwardLevels.lean}{iterate\_phi6\_add\_one}).  Hence

\begin{theorem}[Levels are take-all steps ---
\lean{EM/Population/BackwardLevels.lean}{mem\_preZero\_iff\_sylvWalk\_reaches\_neg\_one}]
\label{thm:levels-are-steps}
$\walkZ{2}{\ell}{N} + 1 \in \mathrm{PreZero}(\ell)$ if and only if the
take-all walk started at $\walkZ{2}{\ell}{N}$ reaches~$-1$.  The hit level
is the number of take-all steps before death.
\end{theorem}

\noindent
So level~$0$ is ``the walk value is already $-1$'', level~$1$ is
$q(w) = -1$, i.e.\ $\Phi_3(w) = 0$---\emph{the classical condition}---and
level~$k$ is death after $k$ steps.  Two structural facts follow at once.
The levels are \textbf{disjoint}, because $q(-1) = 0$ and $q(0) = 0$, so
an orbit that dies is absorbed at~$0$ and never returns
(\lean{EM/Population/BackwardLevels.lean}{death\_level\_unique}).  And
raising the level buys \textbf{no new moduli}: a hit at any level $\geq 1$
produces a root of $\Phi_3$, hence a primitive cube root of unity, hence
$6 \mid \ell-1$ (again for $\ell \ne 2, 3$;
\lean{EM/Population/BackwardLevels.lean}{six\_dvd\_sub\_one\_of\_death}).
What it buys is a strictly larger target \emph{inside} those moduli.

That target is non-empty, and one ring identity governs every level of
it.  The two $q$-preimages of a point are $y$ and $-1-y$, since
$q(-1-y) = q(y)$; passing from a level-$m$ point $z$ to level $m+1$ means
solving $y^2+y=z$, so it is possible exactly when the discriminant $1+4z$
is a square
(\lean{EM/Population/BackwardLevels.lean}{sylvWalk\_step\_of\_isSquare}),
and the two discriminants produced at the next level satisfy
\begin{equation}
\label{eq:delta-identity}
  (1+4y)\,\bigl(1+4(-1-y)\bigr) \;=\; -3 - 16\,q(y)
\end{equation}
(\lean{EM/Population/BackwardLevels.lean}{preimage\_pair\_discriminant})---a
polynomial identity, with no appeal to Vieta and no field hypothesis.
Write $\Delta(z) = -3-16z$.  Then \emph{the pair of discriminants one level
above $z$ multiplies to $\Delta(z)$}, and the whole tower is governed by
evaluating $\Delta$ along backward orbits.

\begin{theorem}[The lift ---
\lean{EM/Population/BackwardLevels.lean}{exists\_death\_level\_add\_two}]
\label{thm:level-lift}
If $z$ is at level~$m$, $1+4z$ is a square, and $\Delta(z)$ is a
\emph{non}-square, then level $m+2$ is occupied.
\end{theorem}

\noindent
Both levels beyond $\Phi_3$ are instances.  At $z = -1$ (level zero),
$\Delta(-1) = 13$, and \eqref{eq:delta-identity} reads
$(1+4\omega)(1+4\omega^{2}) = 13$
(\lean{EM/Population/BackwardLevels.lean}{cube\_root\_pair\_product\_eq\_thirteen}).
At $z = \omega$ (level one) the constant is
$\Delta(\omega)\Delta(\omega^{2}) = 217 = 7\cdot 31$
(\lean{EM/Population/BackwardLevels.lean}{delta\_pair\_product\_eq\_217}).

\begin{keyresult}
\textbf{Levels two and three}
(\lean{EM/Population/BackwardLevels.lean}{exists\_death\_level\_two},
 \lean{EM/Population/BackwardLevels.lean}{exists\_death\_level\_three\_of\_split}).
Suppose $\Phi_3$ has a root modulo $\ell$.
\begin{enumerate}[nosep,leftmargin=*]
\item If $13$ is a quadratic non-residue, level two is occupied.
\item If both level-two branches are present and $217$ is a quadratic
  non-residue, level three is occupied.
\end{enumerate}
Both constants are $\equiv 1 \pmod 4$, so the Jacobi reciprocity sign is
trivial and each criterion is a congruence condition on~$\ell$ alone:
$13$ (resp.\ $217$) is a non-residue mod $\ell$ iff
$\bigl(\tfrac{\ell}{13}\bigr) = -1$ (resp.\ $\bigl(\tfrac{\ell}{217}\bigr) = -1$)
(\lean{EM/Population/BackwardLevels.lean}{not\_isSquare\_iff\_jacobiSym}).
\end{keyresult}

\noindent
Because the levels are disjoint, each occupied level is a fresh avoidance
condition on the perpetual-primality branch, none of them a consequence of
the $\Phi_3$ obstruction:
\[
  q^{\,k}\bigl(\walkZ{2}{\ell}{N}\bigr) \neq -1
  \qquad\text{for every } k \text{ and every prime } \ell \leq \Prod{2}(N),
\]
written out at $k = 2$ and $k = 3$ as
$w^4+2w^3+2w^2+w+1 \neq 0$ and
$w^8+4w^7+8w^6+10w^5+9w^4+6w^3+3w^2+w+1 \neq 0$ for $w = \walkZ{2}{\ell}{N}$
(\lean{EM/Population/BackwardLevels.lean}{psi\_two\_ne\_zero\_of\_perpetual},
 \lean{EM/Population/BackwardLevels.lean}{psi\_three\_ne\_zero\_of\_perpetual}).

\paragraph{Where the climb stops, and why that is worth knowing.}
Two limits appear at level three, and they are the real information.

\emph{The criteria stop being rational.}  When $13$ is a non-residue only
one level-two branch exists, and the level-three criterion is
$\Delta(\omega)$ being a non-square---a condition on~$\omega$, not on a
rational constant.  That is precisely where the tower ceases to be decided
by congruences on~$\ell$ and becomes a splitting condition in a larger
field.  The general statement is an arboreal-Chebotarev question about the
iterated preimage tree of $q$, and it is not proved here.

\emph{The tower is finite.}  Levels are disjoint subsets of a finite
field, so only finitely many are occupied
(\lean{EM/Population/BackwardLevels.lean}{realizedLevels\_finite}).
Raising the level is a bounded resource: the depth of the backward-orbit
tree of $-1$ is an invariant of $\ell$, and it is what the heuristic
$|\mathrm{PreZero}(\ell)| \asymp \sqrt{\ell}$ actually measures.

So the accounting is this.  The classical $\Phi_3$ obstruction is level one
of a tower driven by the single identity \eqref{eq:delta-identity}; levels
two and three are occupied on explicit congruence classes of $\ell$,
produced by reciprocity rather than by search; the moduli never widen; the
tower has a top; and the barrier is untouched, because the criterion still
asks where \emph{one specific orbit} sits.  What has changed is that the
obstruction is now a tower with a known engine, and the first two rungs
above the classical one have been climbed.

%% =========================================================================
\subsection{The Arboreal Tower}
\label{sec:arboreal}
%% =========================================================================

\paragraph{The arboreal tower: Chebotarev is not the missing ingredient.}
Climbing rung by rung terminates, so the natural next move is to ask about
the whole tree at once.  That is the arboreal picture, where the tool of
record is the Chebotarev density theorem applied to the splitting fields of
the iterates: for how many primes $\ell$ is level~$n$ occupied?  Carrying
it out reverses the expected verdict.

Write $\Psi_n(w) = q^{\,n}(w)+1$, so that level~$n$ is occupied at~$w$
exactly when $\Psi_n(w) = 0$.  Two identities organise the tree.  First
$\Psi_n(0) = 1$, because $q(0)=0$: \emph{every level polynomial has
constant term one}.  Second, the tree carries the Euclid--Mullin
accumulator identity itself,
\begin{equation}
\label{eq:tree-accumulator}
  q^{\,n+1}(w) \;=\; w \prod_{j \leq n} \Psi_j(w)
\end{equation}
(\lean{EM/Population/ArborealTower.lean}{sylvWalk\_iterate\_succ\_eq\_prod}),
whence the level values are pairwise coprime
(\lean{EM/Population/ArborealTower.lean}{coprime\_level\_values}) --- the
Euclid cascade of \S\ref{sec:bag}, one level up.

The first identity settles the qualitative half of the Chebotarev question
without any density theorem.  Feeding $q^{\,n}$ a multiple of $B!$ returns
a value $\equiv 1 \pmod{B!}$, so every prime factor of it exceeds~$B$.

\begin{theorem}[Arboreal existence, unconditional ---
\lean{EM/Population/ArborealTower.lean}{exists\_large\_prime\_level\_occupied}]
\label{thm:arboreal-existence}
For every level~$n$ and every bound~$B$ there is a prime $\ell > B$ at which
level~$n$ is occupied; hence infinitely many, for every level
(\lean{EM/Population/ArborealTower.lean}{levelPrimes\_infinite}).
\end{theorem}

\noindent
This is Euclid's argument, in a paper about Euclid's argument.  Chebotarev
would upgrade ``infinitely many'' to ``a positive density of''---and that
upgrade is not what is missing, as the same construction shows when it is
pointed at the orbit rather than at an auxiliary multiple of $B!$.

Take $w = \Prod{2}(N)$.  A prime $\ell$ puts the walk value
$\walkZ{2}{\ell}{N}$ at level~$k$ exactly when $\ell \mid \Psi_k(\Prod{2}(N))$
(\lean{EM/Population/ArborealTower.lean}{level\_witness\_iff}), and
$\Psi_k(\Prod{2}(N))$ is the $k$-th Sylvester tower term above the Euclid
number (\lean{EM/Population/ArborealTower.lean}{tower\_eq\_sylvNat}).  So a
witness prime always exists---take the least prime factor---and by
coprimality the witnesses at distinct levels are \emph{distinct} primes
(\lean{EM/Population/ArborealTower.lean}{minFac\_level\_injective}).  Every
stage~$N$ thus carries an infinite sequence of distinct primes, free of
charge.

\begin{keyresult}
\textbf{Witnesses are free; smallness is everything}
(\lean{EM/Population/ArborealTower.lean}{infinitelyManyComposite\_iff\_witness\_proper}).
\[
  (\mathrm{C}\infty)
  \iff
  \text{for every } N \text{ some witness } \minFac\bigl(\Psi_k(\Prod{2}(N))\bigr)
  \text{ is a \emph{proper} factor.}
\]
On the perpetual-primality branch the witness at level~$k$ is forced to be
$\Psi_k(\Prod{2}(N))$ itself---a prime larger than $\Prod{2}(N)$---which is
exactly the branch
(\lean{EM/Population/ArborealTower.lean}{witness\_eq\_self\_of\_perpetual}).
\end{keyresult}

\noindent
That is the whole of it, and it is the same shape as $(\mathrm{S})$ in
\S\ref{sec:composite-floor}: not \emph{does something exist}, but \emph{is
the thing that exists small}.  A density theorem over $\ell$ cannot supply
it, because the prime it must produce has to divide one specific integer.
The arboreal picture therefore makes Dead End~\#90 as sharp as it can be
made: \textbf{the arboreal input is free, and the residual gap is disjoint
from it.}  Chebotarev is not the missing ingredient, and neither is any
strengthening of it.


%% =========================================================================
\subsection{The Same Telescope over $\FF_p[t]$}
\label{sec:ff-telescope}
%% =========================================================================

The function-field analogue of \S\ref{sec:function-field} is the one
setting in which every analytic input is a theorem rather than a
hypothesis, so it is the natural place to test whether the growth axis is
obstructed by analysis or by something else.  The answer is unambiguous,
and it is the sharpest evidence in this paper that the obstruction was
never analytic.

The growth axis of \S\ref{sec:defect-telescope} transfers to $\Fpt$, and
this is the one place where the function-field model is unambiguously
better behaved.  Over $\ZZ$ the defect telescope carries two
approximations, both artefacts of taking logarithms of integers: the
defect $\delta(n)$ can be slightly negative, so the normalised sequence is
only antitone after a summable correction; and $\delta(n)$ is a real
number, so nothing separates ``the selected factor is proper'' from ``the
selected factor is almost everything''.

Over $\Fpt$ both vanish.  Degrees are additive exactly, and
$\deg(\ffprod(n)+1) = \deg\ffprod(n)$ on the nose
(Lemma~\ref{lem:ff-degree-growth} and
\lean{EM/FunctionField/Analog.lean}{ffProdPlusOne\_natDegree}), so with
$D(n) = \deg \ffprod(n)$ and
$\delta_{\mathrm{FF}}(n) = D(n) - \deg \ffseq(n+1)$ one has the exact
integer recursion
\[
  D(n+1) + \delta_{\mathrm{FF}}(n) \;=\; 2D(n),
  \qquad \delta_{\mathrm{FF}}(n) \in \{0,1,\dots,D(n)-1\}
\]
(\lean{EM/FunctionField/DegreeTelescope.lean}{ffDeg\_succ\_add\_ffDefect}),
telescoping---since $D(0) = 1$---to
\[
  \frac{D(N)}{2^{N}} \;=\; 1 - \sum_{n<N}\frac{\delta_{\mathrm{FF}}(n)}{2^{\,n+1}}
\]
(\lean{EM/FunctionField/DegreeTelescope.lean}{ffNormDeg\_eq\_sub\_sum}).
The right-hand side is antitone with \emph{no} correction term
(\lean{EM/FunctionField/DegreeTelescope.lean}{ffNormDeg\_antitone}), so
$C_{\mathrm{FF}} = \lim_N D(N)/2^{N}$ exists in $[0,1]$, and
$\delta_{\mathrm{FF}}(n) = 0$ holds \emph{if and only if} $\ffprod(n)+1$
is irreducible
(\lean{EM/FunctionField/DegreeTelescope.lean}{ffDefect\_eq\_zero\_iff})---an
equivalence, where the integer statement is only an approximation.  In
particular
\[
  C_{\mathrm{FF}} = 0 \;\Longrightarrow\;
  \ffprod(n)+1 \text{ reducible for infinitely many } n,
\]
which here is a two-line argument---on the perpetually-irreducible branch
$D(N)/2^{N}$ is eventually \emph{constant} and positive---against the
$\log_2\log_2$ chase the integer case requires
(\lean{EM/FunctionField/DegreeTelescope.lean}{ffInfinitelyManyReducible\_of\_ffGrowthConstant\_eq\_zero}).

Trial division transfers too, and here it is an exact inequality on
integers rather than an estimate.  If $\ffprod(n)+1$ is reducible, the
cofactor carries a monic irreducible factor, which minimality forces to be
at least as large as the selected one; hence
$2\deg \ffseq(n+1) \leq D(n)$, i.e.\ $\delta_{\mathrm{FF}}(n) \geq D(n)/2$
(\lean{EM/FunctionField/DegreeTelescope.lean}{two\_mul\_ffDeg'\_le\_of\_not\_irreducible}).
One reducible stage therefore multiplies $D(N)/2^N$ by exactly $3/4$---no
correction term, unlike over $\ZZ$---giving the sharp
\[
  \deg \ffprod(N)
  \;\leq\;
  \Bigl(\tfrac34\Bigr)^{\#\{n<N\,:\,\ffprod(n)+1 \text{ reducible}\}}\, 2^{N}
\]
(\lean{EM/FunctionField/DegreeTelescope.lean}{ffDeg\_le\_pow\_mul}), and
$C_{\mathrm{FF}}$ is a \emph{complete invariant}:
$C_{\mathrm{FF}} = 0$ iff $\ffprod(n)+1$ is reducible infinitely often,
and $C_{\mathrm{FF}} > 0$ iff the sequence is eventually perpetually
irreducible
(\lean{EM/FunctionField/DegreeTelescope.lean}{ffInfinitelyManyReducible\_iff\_ffGrowthConstant\_eq\_zero},
 \lean{EM/FunctionField/DegreeTelescope.lean}{ffGrowthConstant\_pos\_iff\_perpetual}).
This is Theorem~\ref{thm:defect-dichotomy} over $\Fpt$, with every epsilon
removed.

And yet the question does not move---which is the sharpest available form
of the barrier thesis.  Both settings reduce $(\mathrm{C}\infty)$ to a
single real number defined \emph{from the orbit}, and in both that number
is beyond any finite computation, since a defect at stage~$n$ is
discounted by $2^{-n}$.  What would settle it is the least-degree
irreducible factor of the \emph{one} polynomial $\ffprod(n)+1$; the density
of irreducibles of each degree, which over $\Fpt$ is not merely known but
exactly computable, is a statement about the population.  This is Dead
End~\#90 in the setting where every analytic input is a theorem---which is
the strongest available evidence that the analytic input was never the
obstruction.

\paragraph{Where the question does move.}
The preceding paragraph is about the question \emph{as posed uniformly
in~$p$}.  For particular primes the model decides it, in both directions.
Over $\FF_5[t]$ the seed-$t$ sequence is perpetually irreducible
(\lean{EM/FunctionField/StableTower.lean}{tower\_euclid\_irreducible};
Theorem~\ref{thm:ffmc-false-five}): $\delta_{\mathrm{FF}}(n) = 0$ for every
$n$, $D(N) = 2^N$ exactly, $C_{\mathrm{FF}} = 1$, and
$(\mathrm{C}\infty)$ is \emph{false}.  The branch that
\S\ref{sec:cinfty-verdict} cannot rule out over~$\ZZ$ is realised there,
and with it the failure of the conjecture.  In the other direction the
floor is a theorem for $p = 2$, where $P \mapsto P(P+1) = P^2 + P$ is
additive and $(F+1)^3P + 1 = (P^4+P^3+1)(P^4+P^3+P^2+P+1)$ forces a
reducible stage among any four consecutive ones (sharp: the seed~$t$
attains three), and for $p \equiv 1 \pmod 3$, where $\Phi_3$ splits and no
two consecutive Euclid polynomials are irreducible.  Which case a prime
$p \equiv 2 \pmod 3$ falls into is decided by a finite check on the orbit
of $-1/4$ under $y \mapsto y^2+y$ in~$\Fp$ (the exceptional primes of
\S\ref{sec:function-field}).  All three statements are machine-checked,
for every FF-EM data:
\lean{EM/FunctionField/CompositeFloors.lean}{ffGrowthConstant\_eq\_zero\_of\_two}
and
\lean{EM/FunctionField/CompositeFloors.lean}{not\_four\_consecutive\_irreducible}
for $p = 2$,
\lean{EM/FunctionField/CompositeFloors.lean}{ffGrowthConstant\_eq\_zero\_of\_one\_mod\_three}
for $p \equiv 1 \pmod 3$, and the $\FF_5$ refutation above.  Over $\FF_3$,
where $\Phi_3 = (y-1)^2$ has a double root, an irreducible Euclid
polynomial is followed by a perfect square and the floor again holds with
constant~$1$ (\lean{EM/FunctionField/CharThree.lean}{ffGrowthConstant\_eq\_zero});
the sharpness of the constant~$3$ over $\FF_2$ is machine-checked too
(\lean{EM/FunctionField/CharTwo.lean}{ff\_two\_attains\_three}).  The
function-field model thus decides $(\mathrm{C}\infty)$ in both directions,
prime by prime, while over~$\ZZ$ it decides nothing.

%% =========================================================================
\subsection{The Verdict}
\label{sec:cinfty-verdict}
%% =========================================================================

Collecting the section: $(\mathrm{C}\infty)$ is required by every route to
$\MC$; it is equivalent to the vanishing of one real number; that number
vanishes if and only if the sequence does \emph{not} eventually become a
Sylvester tower of primes; and the failure branch is exactly the autonomous
branch, on which every obstruction this paper owns applies.  What remains is
a single statement, and it is worth writing it in the plainest form
available.

\begin{keyresult}
\textbf{$(\mathrm{C}\infty)$, plainly.}
There is no $N$ with $\Prod{2}(n)+1$ prime for every $n \geq N$.
Equivalently: no Sylvester tower $s \mapsto s^2-s+1$ seeded at a Euclid
number is all-prime.  Equivalently:
$C = \lim_N \log\Prod{2}(N)/2^{N}$ vanishes.
\end{keyresult}

\noindent
Since $\MC \Rightarrow (\mathrm{C}\infty)$
(\lean{EM/Population/CompositeFloor.lean}{infinitelyManyComposite\_of\_mullin};
a short implication---perpetual primality would leave only finitely many
primes selected), Mullin's conjecture \textbf{implies a
no-infinite-prime-tower statement for a doubly exponential sequence},
which is open.  Every conjecture is at least as hard as its open
consequences, so this is not a lower bound on difficulty in any
technical sense; it is a statement of \emph{shape}: any proof of $\MC$
must, along the way, rule out a Sylvester-type prime tower.  The shape can
be made a little sharper.  The level polynomials $g_k = \Phi_6^{\,k}+1$ are
irreducible over~$\mathbb{Q}$ for every~$k$
(Theorem~\ref{thm:level-irreducible-Q}, proved by reduction mod~$5$), so
perpetual primality from stage~$N$ says that the single integer
$\Prod{2}(N)$ is a simultaneous prime value of infinitely many irreducible
polynomials.  For comparison, even the Bunyakovsky instance ``$y^2+y+1$ is
prime infinitely often'' is open.

\paragraph{The Fermat test.}
The comparison is not loose.  The Fermat numbers $F_n = 2^{2^n}+1$ have the
same shape: doubly exponential, $\sum_n 1/\log F_n < \infty$, believed
composite for all large~$n$---and whether infinitely many $F_n$ are
composite is \emph{open}.  Our $C$ plays the role of $\log 2$ there.
The analogy is one of shape, not of reduction: nothing here relates the
two questions formally, and a method exploiting the recursion
$E_{n+1} = \Phi_3(E_n - 1)$ specific to Euclid numbers need not say
anything about $2^{2^n}+1$.

This suggests a filter, and we have applied it to every technique proposed
in this section:

\begin{center}
\emph{If a proposed method would not settle the Fermat question, it will
not settle $(\mathrm{C}\infty)$.}
\end{center}

\noindent
Congruence and reciprocity fail it (\S\ref{sec:sylvester-tower}); so do
$abc$ and Baker's method, for reasons recorded in
\S\ref{app:dead-ends}; the growth axis fails it because $C$ is a complete
invariant (\S\ref{sec:defect-telescope}); level-by-level climbing fails it
because the ladder is finite (\S\ref{sec:level-tower}); and the arboreal
route fails it because its input is free and its output is a density
statement, where a size condition on one integer is wanted
(\S\ref{sec:arboreal}).  The one method that has ever worked on a question
of this shape---varying the seed, as for generalised Fermat numbers
$a^{2^n}+1$, where the composite-infinitely-often statement \emph{is} known
for almost all~$a$---is available here too, through the ensemble framework
of \S\ref{sec:ensemble-variant}; and it delivers almost-all statements
while leaving the specific orbit untouched.  That is the orbit-specificity
barrier, arrived at for the sixth time in this section and by six
independent routes.

\begin{frontierbox}
\textbf{What would settle it.}
A method that decides, for one explicitly given integer of doubly
exponential size, whether its least prime factor is proper---without a
factorisation and without a density hypothesis.  No such method is known
for any sequence of this shape.  Until one exists, $(\mathrm{C}\infty)$ is
where every route in this paper stops, and the results of
\S\S\ref{sec:growth-floor}--\ref{sec:arboreal} are best read as a map of
that stopping point rather than as progress towards passing it.
\end{frontierbox}
